\section{Введение и основные определения.}
\subsection{Классическое и геометрическое определние вероятности.}
\subsubsection{Задача о разделе ставки (постановки Пачолли и Ферма).}
\subsubsection{Парадокс Монти-Холла.}
\subsubsection{Парадокс Бертрана.}
\subsection{Аксиоматическое определение вероятности.}
\subsubsection{Вероятностное пространство, основные свойства.}
\subsubsection{Минимальная сигма-алгебра, основные свойства. Борелевская сигма-алгебра.}
\subsubsection{Построение вероятностного пространства на примере простейшего потока событий.}
\subsection{Условная вероятность.}
\subsubsection{Независимые события, соотношения между независимостью в совокупности и попарной независимостью.}
\subsubsection{Независимые алгебры и сигма-алгебры.}
\subsubsection{Теорема о <<продолжении>> отношения независимости с алгебры на минимальную сигма-алгебру.}
\subsection{Теорема непрерывности вероятности}
\subsubsection{Леммы Бореля-Кантелли.}
\subsubsection{Формулы полной вероятности и Байеса.}
\subsubsection{Формула включений-исключений для системы событий.}
\subsection{Симметрическая разность и её свойства.}
\subsubsection{Теорема об аппроксимации.}
\subsubsection{Хвостовые алгебры и остаточные события.}
\subsubsection{Закон <<0-1>> Колмогорова для остаточных событий.}
\subsection{Случайные величины.}
\subsubsection{Определение случайной величины, измеримые и борелевские функции.}
\subsubsection{Измеримость $\sup$ и $\inf$, суммы и произведения случайных величин.}
\subsubsection{Критерий существований функциональной зависимости случайных величин.}
\subsection{Функции распределения.}
\subsubsection{Определения и свойства. Дискретные и непрерывные случайные величины.}
\subsubsection{Функции распределения дискретных, абсолютно непрерывных и сингулярных случайных величин.}
\subsubsection{Функция плотности вероятности и её свойства.}
\subsubsection{Условная функция распределения и её свойства.}
\subsubsection{Функция распределения смеси.}
\subsection{Случайный вектор.}
\subsubsection{Функция распределения случайного вектора и её свойства.}
\subsubsection{Понятие независимости случайных величин, функция распределения независимых случайных величин.}